\begin{figure}[htbp]
\begin{syntax}{Type}{A,B}
\sitem{\tpdiam}{diamond}
\sitem{\tpunit}{unit}
\sitem{\tpsum{A}{B}}{sum}
\sitem{\tptens{A}{B}}{tensor}
\sitem{\tparr{A}{B}}{arrow}
\sitem{\tplist{A}}{list}
\\
\newsort{Term}{M,N}
\sitem{\tmnull}{unit intro}
\sitem{\tminj{i}{M} \quad\quad\text{($i \in \{1,2\}$)}}{sum intro}
\sitem{\tmcase{M}{x_1}{N_1}{x_2}{N_2}}{sum elim}
\sitem{\tmpair{M_1}{M_2}}{tensor intro}
\sitem{\tmletp{M}{x_1}{x_2}{N}}{tensor elim}
\sitem{\tmlam{x}{M}}{arrow intro}
\sitem{\tmapp{M}{N}}{arrow elim}
\sitem{\tmnil}{list intro}
\sitem{\tmcons{M_0}{M_1}{M_2}}{list intro}
\sitem{\tmrec{M}{N_1}{x_0}{x_1}{x_2}{N_2}}{list elim}
\end{syntax}
\Description{An inductive definition of the syntax for LFPL types and terms}
\caption{LFPL's type and term syntax}
\label{fig:tp-tm-syntax}
\end{figure}